% ---------------------------------------------------------------------------
\section{Reduced spherical dynamics of a fully activated PFC droplet}\label{sec:spherical}
The ceilings of Section~\ref{sec:thermo} say how much work is available; this section computes how a
spherical cavity actually spends it, and which unknown controls the collapse. Spherical symmetry
cannot produce a directed jet (Section~\ref{sec:jet}); it gives the expansion, the collapse
violence and the content pressure history that the jet calculation needs.

\subsection{Step S1 --- Carrier: Keller--Miksis equation (sourced)}
\emph{Symbols.} $R(t)$, $\dot R$, $\ddot R$ wall radius, velocity, acceleration; $\rho_{\mathrm w}$,
$c_{\mathrm w}$, $\mu$ carrier density, sound speed and viscosity (for a gel, $\mu=\mu_{\mathrm w}+\mu_{\mathrm{gel}}$);
$\sigma$ wall tension; $p_b$ cavity pressure; $S_{\mathrm{el}}(R)$ elastic stress integral of the gel
($S_{\mathrm{el}}=0$ in water);
$p_L$ carrier pressure at the wall. The liquid is weakly compressible (outgoing spherical
waves, first order in $\dot R/c_{\mathrm w}$; Keller \& Miksis 1980). Before the two relations:
$R,a_0$ are lengths (m), $t$ time (s), overdot time differentiation, $\rho_{\mathrm w}$ density
(kg/m$^3$), $c_{\mathrm w}$ speed (m/s), $\mu$ viscosity (Pa\,s), $\sigma$ tension (N/m),
$p_b,p_L,\pinf,S_{\mathrm{el}}$ pressures (Pa), and $G$ shear modulus (Pa):
\begin{gather}
\Bigl(1-\frac{\dot R}{c_{\mathrm w}}\Bigr)R\ddot R+\frac32\Bigl(1-\frac{\dot R}{3c_{\mathrm w}}\Bigr)\dot R^2
=\Bigl(1+\frac{\dot R}{c_{\mathrm w}}\Bigr)\frac{p_L-\pinf}{\rho_{\mathrm w}}+\frac{R}{\rho_{\mathrm w}c_{\mathrm w}}\frac{\dd p_L}{\dd t},
\label{eq:km}\\
p_L=p_b-\frac{2\sigma}{R}-\frac{4\mu\dot R}{R}+S_{\mathrm{el}}(R),\qquad
S_{\mathrm{el}}(R)=-\frac G2\Bigl[5-4\frac{a_0}{R}-\Bigl(\frac{a_0}{R}\Bigr)^4\Bigr].\label{eq:gelS}
\end{gather}
The gel term is the neo-Hookean stress integral (identical to Finite-PFC E08 with
$S_{\mathrm{el}}=-p_{\mathrm{el}}$; stress-free at the initial droplet radius $a_0$). Because $\dd p_L/\dd t$ contains
$-4\mu\ddot R/R$, that term is moved to the left. For the explicit equation, $R>0$ (m),
$\dot R$ (m/s), $\ddot R$ (m/s$^2$), $\rho_{\mathrm w}>0$ (kg/m$^3$), $c_{\mathrm w}>0$ (m/s),
$p_L,\pinf,p_b,S_{\mathrm{el}}$ (Pa), $\dot p_b$ (Pa/s), $\sigma$ (N/m), $\mu$ (Pa\,s), and $S'_{\mathrm{el}}(R)$
(Pa/m) have the meanings above. The denominator is positive while $|\dot R|/c_{\mathrm w}<1$:
\begin{equation}
\ddot R=\frac{\bigl(1+\tfrac{\dot R}{c_{\mathrm w}}\bigr)\tfrac{p_L-\pinf}{\rho_{\mathrm w}}
-\tfrac32\bigl(1-\tfrac{\dot R}{3c_{\mathrm w}}\bigr)\dot R^2+\tfrac{R}{\rho_{\mathrm w}c_{\mathrm w}}
\bigl(\dot p_b+\tfrac{2\sigma\dot R}{R^2}+\tfrac{4\mu\dot R^2}{R^2}+S'_{\mathrm{el}}(R)\dot R\bigr)}
{\bigl(1-\tfrac{\dot R}{c_{\mathrm w}}\bigr)R+\tfrac{4\mu}{\rho_{\mathrm w}c_{\mathrm w}}} .\label{eq:km_explicit}
\end{equation}
\emph{Check:} for $c_{\mathrm w}\to\infty$, $\mu=\sigma=G=0$, $p_b=0$, \eqref{eq:km_explicit} reduces to
the Rayleigh equation; the implementation reproduces the Rayleigh collapse time
$0.914681\,R_{\max}\sqrt{\rho/\Delta p}$ to $1.7\times10^{-8}$ and Limits E17 to $2\times10^{-10}$
(Appendix~\ref{app:verif}). \emph{Validity:} first order in Mach number; runs are stopped when
$|\dot R|=0.3\,c_{\mathrm w}$ (flag ``M''): beyond that the peak collapse pressure is not resolved
by this model and is reported as unresolved.

\paragraph{Why the ideal singular endpoint is not a maximum.}
In the incompressible, inviscid, tension-free limit with a constant cavity pressure,
let $\Delta p=\pinf-p_b>0$ (Pa), $\rho_{\mathrm w}>0$ (kg/m$^3$), and $\Rm>0$ (m).
The wall starts at $R=\Rm$, $\dot R=0$ and collapses, so $\dot R\le0$; all radii are in m,
time in s. Define $w(R)=\dot R^2$ (m$^2$/s$^2$). The Rayleigh limit of
Eq.~\eqref{eq:km_explicit} (Brennen, Chapter 2) can be integrated without using its singular endpoint:
\begin{align}
R\ddot R+\tfrac32\dot R^2&=-\frac{\Delta p}{\rho_{\mathrm w}},\qquad
\ddot R=\tfrac12\frac{\dd w}{\dd R},\notag\\
\frac{\dd}{\dd R}(R^3w)&=-\frac{2\Delta p}{\rho_{\mathrm w}}R^2,\notag\\
R^3w-\Rm^3(0)&=-\frac{2\Delta p}{3\rho_{\mathrm w}}(R^3-\Rm^3),\notag\\
\dot R&=-\sqrt{\frac{2\Delta p}{3\rho_{\mathrm w}}
\left[\left(\frac{\Rm}{R}\right)^3-1\right]}\quad(0<R\le\Rm).
\label{eq:rayleigh_speed}
\end{align}
The first substitution uses $\ddot R=\dot R\,\dd\dot R/\dd R$ away from the initial turning point;
the integrated solution extends continuously to it. Thus the ideal wall speed diverges as
$R^{-3/2}$ when $R\to0$. This is compatible with finite energy: with the radial liquid velocity
$u_r(r)=R^2\dot R/r^2$ (m/s), liquid radius $r\ge R$ (m), and radial kinetic energy $K_r$ (J),
\begin{equation}
K_r=\int_R^\infty\tfrac12\rho_{\mathrm w}u_r^2\,4\pi r^2\dd r
=2\pi\rho_{\mathrm w}R^3\dot R^2
=\tfrac{4\pi}{3}\Delta p(\Rm^3-R^3).
\label{eq:rayleigh_energy}
\end{equation}
The increasing velocity concentrates in a shrinking region. A point speed can therefore diverge
without violating a finite budget; a prescribed participating mass has the separate finite bound
of Eq.~\eqref{eq:Kj}. Compressibility, vapour content, viscosity and asymmetric focusing intervene
before $R=0$. The ideal singular endpoint is a model breakdown, not a physical laser-cavitation
maximum or a PFC jet-speed prediction.

\subsection{Step S2 --- Content during expansion: equilibrium isentrope of the real fluid}
At the end of heating the droplet is metastable liquid at $(\Tk,\pinf)$ with specific internal energy
$u_{\mathrm a}$ and specific volume $v_{\mathrm a}$. It first relaxes at fixed $(u_{\mathrm a},v_{\mathrm a})$ to the
equilibrium two-phase state (pressure $p_{\mathrm{relax}}$, entropy $s_{\mathrm{relax}}$; this step is
irreversible). It then expands along the equilibrium isentrope $s=s_{\mathrm{relax}}$.
Here $p_b,p_{\mathrm{eq}}$ are pressures (Pa), $T_b,T_{\mathrm{eq}}$ temperatures (K),
$s_{\mathrm{relax}}$ fixed specific entropy (J/(kg\,K)), $V$ volume (m$^3$), $\mF>0$ PFC mass
(kg), $v_{\mathrm a}$ initial specific volume (m$^3$/kg), and $0$ the initial time:
\begin{equation}
p_b=p_{\mathrm{eq}}(s_{\mathrm{relax}},V/\mF),\qquad T_b=T_{\mathrm{eq}}(s_{\mathrm{relax}},V/\mF),
\qquad V(0)=\mF v_{\mathrm a}.\label{eq:isentrope}
\end{equation}
For the 6\,$\mu$m benchmark $p_{\mathrm{relax}}=\PFPkmprelax$\,MPa (PFP) and $\PFHkmprelax$\,MPa (PFH).
\emph{Justification of adiabaticity:} the thermal diffusion length in water over the expansion time
($\approx1\,\mu$s) is $\sqrt{\alpha_{\mathrm w}t}\approx0.4\,\mu$m $\ll R$; the content is nearly uniform
because its sound speed ($\sim$100\,m/s) exceeds the wall speed. Phase 1 ends when the PFC is
entirely vapour (quality 1).

\subsection{Step S3 --- Content after full vaporisation: vapour, heat and mass transfer}
\emph{Symbols.} $m_{\mathrm v}$ PFC vapour mass (the condensed film has mass $\mF-m_{\mathrm v}$); $m_{\mathrm a}$
air mass; $T_b$ uniform core temperature; $c_{v\mathrm v},c_{p\mathrm v}$ ideal-gas heat capacities of the
vapour from the EOS ($k_{\mathrm v}=c_{p\mathrm v}/c_{v\mathrm v}$); $\mathcal R_{\mathrm v}=\mathcal R_u/M$;
$E_b=(m_{\mathrm v}c_{v\mathrm v}+m_{\mathrm a}c_{v\mathrm a})T_b$ content internal energy (ideal-gas reference
$u(0)=0$); $T_i$ wall temperature; $j$ net evaporative mass flux (kg\,m$^{-2}$s$^{-1}$); $q_{\mathrm g}$
conductive heat flux from the core to the wall; $A=4\pi R^2$.
For all three rows, $m_{\mathrm v},m_{\mathrm a},E_b$ are respectively masses (kg) and energy (J),
$\mathcal R_{\mathrm v},\mathcal R_{\mathrm a}$ specific gas constants (J/(kg\,K)),
$T_b,T_i>0$ temperatures (K), $V>0$ volume (m$^3$), $A$ area (m$^2$), $R$ radius (m),
$p_b,p_{\mathrm v},\psat$ pressures (Pa), $c_{p\mathrm v}$ heat capacity (J/(kg\,K)),
$k_{\mathrm g}$ vapour conductivity (W/(m\,K)), $q_{\mathrm g}$ flux (W/m$^2$),
$j$ mass flux (kg/(m$^2$\,s)), and $\alpha_{\mathrm c}\in[0,1]$ dimensionless coefficient.
Overdots are time derivatives; $\pi$ is the dimensionless circle constant:
\begin{align}
p_b&=\frac{(m_{\mathrm v}\mathcal R_{\mathrm v}+m_{\mathrm a}\mathcal R_{\mathrm a})T_b}{V},\qquad
\dot m_{\mathrm v}=A\,j,\label{eq:content}\\
\dot E_b&=-p_b\dot V+A\,j\,c_{p\mathrm v}T_i-A\,q_{\mathrm g},\qquad
q_{\mathrm g}=k_{\mathrm g}\frac{T_b-T_i}{R/\pi},\label{eq:km_energy}\\
j&=\frac{2\alpha_{\mathrm c}}{2-\alpha_{\mathrm c}}\frac{1}{\sqrt{2\pi\mathcal R_{\mathrm v}}}
\Bigl[\frac{\psat(T_i)}{\sqrt{T_i}}-\frac{p_{\mathrm v}}{\sqrt{T_b}}\Bigr],\qquad p_{\mathrm v}=\frac{m_{\mathrm v}\mathcal R_{\mathrm v}T_b}{V}.\label{eq:hks}
\end{align}
\eqref{eq:km_energy} is the first law of the open content (Limits E29) with the vapour entering at the
wall temperature; \eqref{eq:hks} is the Hertz--Knudsen--Schrage law. The gas-side conduction uses the
boundary-layer length $R/\pi$ of Toegel et al.\ (2000). Evaporation is switched off smoothly when the
condensed inventory is exhausted (finite PFC inventory, Limits E24); condensation is always allowed.
The switch from the real EOS to the ideal vapour is pressure-continuous; the implied temperature
changes from $\PFPkmTswitch$\,K (EOS) to $\PFPkmTswitchIdeal$\,K (ideal) for PFP, the non-ideality
($Z\approx0.96$) of the vapour at that state. This pressure-continuous switch does not establish
energy or entropy continuity of the two reference conventions. Subsequent runs with
$p_b>0.3p_{\mathrm c}$ are flagged as an ideal-vapour warning: a numerical minimum can still be
an EOS extrapolation. The flag is not a validation certificate.

\subsection{Step S4 --- Wall energy balance and carrier heat diffusion}
The carrier temperature $T(y,t)$ is written in the Lagrangian volume coordinate $y=(r^3-R^3)/3$.
\emph{Derivation:} with $u_r$ radial carrier velocity (m/s), incompressibility gives
$r^2u_r=R^2\dot R$, so $\dd y/\dd t=r^2u_r-R^2\dot R=0$ for a
material point ($y$ is a material label); with $\partial_r=r^2\partial_y$ the advection--diffusion equation
$\partial_tT+u_r\partial_rT=\alpha_{\mathrm w}r^{-2}\partial_r(r^2\partial_rT)$ becomes the next
equation. Here $T$ is temperature (K), $t$ time (s), $y$ material volume coordinate (m$^3$),
$r,R$ radii (m), $\alpha_{\mathrm w}>0$ thermal diffusivity (m$^2$/s), and
$\partial_t|_y$ the derivative at fixed material label:
\begin{equation}
\partial_tT\big|_y=\alpha_{\mathrm w}\,\partial_y\bigl(r^4\partial_yT\bigr),\qquad r=(3y+R^3)^{1/3},\label{eq:heat}
\end{equation}
with $T(y\to\infty)=T_0$. At the wall, $k_{\mathrm w}$ is conductivity (W/(m\,K)),
$R$ radius (m), $T(y,t),T_i$ temperatures (K), $y$ volume coordinate (m$^3$),
$j$ mass flux (kg/(m$^2$\,s)), $h_{\mathrm{fg}}(T_i)$ latent heat (J/kg), and $q_{\mathrm g}$ gas-side heat flux
(W/m$^2$). The interfacial energy balance is
\begin{equation}
k_{\mathrm w}\,R^2\,\partial_yT\big|_{y=0}=j\,h_{\mathrm{fg}}(T_i)-q_{\mathrm g},\label{eq:wallbalance}
\end{equation}
solved for $T_i$ by Newton iteration. $h_{\mathrm{fg}}(T)$ is the EOS latent heat. Equation \eqref{eq:heat} is
discretised by finite volumes on a geometrically stretched grid (80 nodes; convergence in
Appendix~\ref{app:verif}).

\subsection{Step S5 --- Energy ledger}
Multiplying \eqref{eq:km} (incompressible part) by $4\pi R^2\rho_{\mathrm w}\dot R$ and integrating from
activation to the spherical maximum radius (where radial kinetic energy vanishes).
Here $p_b,\pinf$ are pressures (Pa), $\dd V,\Delta V$ volume increments (m$^3$),
$\sigma$ tension (N/m), $\Delta A$ area change (m$^2$), $E_{\mathrm{gel}},D_{\mathrm{visc}},
E_{\mathrm{rad}}$ energies (J), $\mu$ viscosity (Pa\,s), $R$ radius (m), $\dot R$ wall speed (m/s),
and $\dd t$ time increment (s); the integrations cover that expansion interval:
\begin{equation}
\int p_b\,\dd V=\pinf\Delta V+\sigma\Delta A+E_{\mathrm{gel}}+D_{\mathrm{visc}}+E_{\mathrm{rad}},\qquad
D_{\mathrm{visc}}=\int16\pi\mu R\dot R^2\dd t ,\label{eq:ledger}
\end{equation}
where $E_{\mathrm{rad}}$ (acoustic radiation, the compressible remainder) is obtained by difference.
For the 6\,$\mu$m PFP droplet in water (frozen closure): heat $Q=\PFPledgerQ$\,nJ; exergy
$\Phia=\PFPledgerPhi$\,nJ; content work $\int p_b\dd V=\PFPledgerW$\,nJ, of which $\PFPledgerpV$\,nJ
displaces the ambient (not useful), $\PFPledgerSurf$\,nJ is surface energy, $\PFPledgerVisc$\,nJ
viscous loss, $\PFPledgerRad$\,nJ radiated. The useful work is $\PFPledgerUseful$\,nJ and the
exergy still held by the cold vapour at $R_{\max}$ is $\PFPledgerContent$\,nJ: the expansion converts
only about one fifth of $\Phia$ into stored liquid/surface energy, and most of the remainder stays
in the vapour as the potential to recondense. The residual-based mechanical bookkeeping closes to
$\PFPledgerClosure$ (relative), but this is not an independent radiation or thermochemical conservation
test: radiation was defined by difference. The full exergy remainder is
$\PFPledgerExergyResidual$\,nJ ($\PFPledgerExergyResidualPercent$\% of $\Phia$) for PFP and
$\PFHledgerExergyResidual$\,nJ ($\PFHledgerExergyResidualPercent$\%) for PFH. The negative PFP
remainder exposes a small reference/closure inconsistency across the adopted phase models; it cannot
represent physical negative exergy destruction. Thus the rigorous availability inequality is a separate
thermodynamic result, while this reduced ledger estimates how the budget is spent without an assumed
efficiency. The independent incompressible frozen mechanical-identity test is listed in
Appendix~\ref{app:verif}.

\subsection{Results: condensation decides the collapse}\label{sec:km_results}
Table~\ref{tab:km} and Fig.~\ref{fig:km} summarise 216 runs (PFP/PFH; $d_0=3,6,12\,\mu$m; water,
5\,kPa and 20\,kPa gels; four closures; $x_{\mathrm{air}}=0,0.3$). Findings:
\begin{enumerate}[leftmargin=*]
\item \textbf{Expansion is robust.} Maximum expansion wall speeds are $18$--$20$\,m/s in water for
every size and closure (PFP 6\,$\mu$m: $\PFPkmUexp$\,m/s, $R_{\max}=\PFPkmRmax\,\mu$m against the
rigorous $R^*=\PFPRboundSix\,\mu$m); the far-field expansion pressure at 50\,$\mu$m is
$\PFPkmpexpFifty$\,MPa. The equilibrium expansion is consistent with the H5-conditioned source
ceiling (Proposition~\ref{prop:psrc}); this does not prove its non-equilibrium extension.
\item \textbf{Collapse is controlled by one unmeasured number, $\alpha_{\mathrm c}$.} For
$\alpha_{\mathrm c}\le0.1$ the vapour cannot condense fast enough and cushions the collapse
(PFP 6\,$\mu$m: $R_{\min}=\PFPkmRminAlphaTenth\,\mu$m, $U_{\mathrm c}=\PFPkmUcolAlphaTenth$\,m/s);
at $\alpha_{\mathrm c}=0.3$, the original 80-node PFP sample gives
$U_{\mathrm c}=\PFPkmUcolAlphaThird$\,m/s (thermal refinement below); the 6\,$\mu$m water cases reach
the Keller--Miksis screen (444.6\,m/s) at $\alpha_{\mathrm c}=0.5$ for PFP and already at
$\alpha_{\mathrm c}=0.3$ for PFH. Their collapse peaks are unresolved. These are sampled,
species-dependent transitions, not exact universal thresholds. The
frozen limit gives $U_{\mathrm c}=\PFPkmUcolFrozen$\,m/s. No measurement of $\alpha_{\mathrm c}$ for PFC
vapour on water was found; it is the single most important missing datum for collapse pressures.
\item \textbf{PFH collapses more violently than PFP at equal $\alpha_{\mathrm c}$,} because its vapour
is more subcooled at $T_0$ ($\psat=\PFHpsatZero$ versus $\PFPpsatZero$\,kPa): condensation
is thermodynamically favoured and the cushion pressure is lower.
\item \textbf{Dissolved air at $x_{\mathrm{air}}\le0.3$ is negligible} (change of $R_{\min}$ $\le\airEffectRmin$\%)
because the air mass is tiny compared with the PFC vapour mass.
\item \textbf{Gel stiffness suppresses or eliminates the collapse.} The expansion of small PFP droplets
arrests while still two-phase (flag ``A''): 3 and 6\,$\mu$m at 20\,kPa (6\,$\mu$m: $R=\gelArrestR\,\mu$m, quality
$\gelArrestQ$) and 3\,$\mu$m at 5\,kPa; the content then follows the reversible equilibrium isentrope and no violent
collapse occurs in this model. PFH droplets, which carry more exergy per mass, still vaporise fully in both gels.
For the 6\,$\mu$m PFP droplet in a 5\,kPa gel the collapse speed stays below 15\,m/s for $\alpha_{\mathrm c}\le1$;
12\,$\mu$m droplets collapse violently for $\alpha_{\mathrm c}\gtrsim0.5$ even in gel.
\end{enumerate}
\begin{figure}[ht]\centering
\includegraphics[width=0.95\textwidth]{figures/km_alpha.pdf}
\caption{Spherical collapse of the fully activated PFC cavity in water versus the condensation
coefficient. Above $\alpha_{\mathrm c}\approx0.3$--$0.5$ the collapse runs into the compressible regime
(dash-dot line); below $\approx0.1$ the vapour cushions it.}\label{fig:km}
\end{figure}
\input{generated/table_km.tex}

\subsection{Thermal-grid refinement and vapour-model warning}\label{sec:km_refine}
The selected 6\,$\mu$m air-free water cases at $\alpha_{\mathrm c}=0.2,0.3,0.5$ were rerun with
$N_{\mathrm{th}}=160,320$ thermal nodes, changing no other parameter from the 80-node records.
Table~\ref{tab:kmrefine} separates uncensored reduced-model peaks from Mach-event values.
For PFP at $\alpha_{\mathrm c}=0.3$, the reported wall speed rises from
$\PFPkmRefineThirdUBase$ to $\PFPkmRefineThirdUMid$ to $\PFPkmRefineThirdUFine$\,m/s:
$\PFPkmRefineThirdDeltaU$\% from 160 to 320 nodes and $\PFPkmRefineThirdDeltaUFull$\% from 80 to 320.
This transitional case does not meet a 2\% speed-change criterion at the finest pair.
The uncensored $\alpha_{\mathrm c}=0.2$ speeds are more stable, but all six refined cases still
trigger the ideal-vapour warning; numerical convergence does not cure that physical closure limit.

For the censored PFP $\alpha_{\mathrm c}=0.5$ run, the reported pre-stop radius, vapour mass,
content pressure and stopping time change by $\PFPkmRefineHalfDeltaR$\%,
$\PFPkmRefineHalfDeltaMass$\%, $\PFPkmRefineHalfDeltaPb$\% and
$\PFPkmRefineHalfDeltaTime$\%, respectively, from 160 to 320 nodes. Equal speeds of 444.6\,m/s
at the screen are imposed by the stopping rule and do not demonstrate peak convergence.
The cushioning/violent-collapse classification is unchanged by this refinement, while the
physical peak pressure and fast-collapse endpoint remain unresolved.
\input{generated/table_km_refine.tex}
